% Part: first-order-logic
% Chapter: beyond
% Section: higher-order-logic

\documentclass[../../../include/open-logic-section]{subfiles}

\begin{document}

\olfileid{fol}{byd}{hol}

\olsection{منطق مرتبهٔ بالاتر}

گذار از منطق مرتبهٔ اول به منطق مرتبهٔ دوم به ما امکان داد درون زبان
صوری دربارهٔ مجموعه‌های اشیای !!{domain} مرتبهٔ اول سخن بگوییم. چرا
همین‌جا متوقف شویم؟ برای نمونه، منطق مرتبهٔ سوم باید به ما امکان دهد با
مجموعه‌های مجموعه‌های اشیا، یا شاید حتی مجموعه‌هایی شامل هم اشیا و هم
مجموعه‌های اشیا، سروکار داشته باشیم. منطق مرتبهٔ چهارم نیز اجازه خواهد
داد دربارهٔ مجموعه‌های اشیایی از این نوع سخن بگوییم. چنان‌که شاید حدس
زده باشید، این ایده را می‌توان به‌دلخواه تکرار کرد.

در عمل، منطق مرتبهٔ بالاتر غالباً بر حسب تابع‌ها !!{formula}‌بندی
می‌شود، نه رابطه‌ها. (با صرف نظر از یکی‌گرفتن‌های طبیعی، این تفاوت
اساسی نیست.) با داشتن چند «نوع» پایهٔ $A$، $B$، $C$،~\dots (که از این
پس آن‌ها را «تیپ» می‌نامیم)، می‌توانیم با وضع قاعدهٔ زیر تیپ‌های تازه‌ای
بسازیم:
\begin{quote}
اگر $\sigma$ و~$\tau$ تیپ‌های متناهی باشند، آنگاه $\sigma \to \tau$
نیز تیپی متناهی است.
\end{quote}
تیپ‌ها را «برچسب»های نحوی‌ای در نظر بگیرید که اشیای مورد نظر ما در
!!{domain} را طبقه‌بندی می‌کنند؛ $\sigma \to \tau$ آن اشیایی را توصیف
می‌کند که تابع‌هایی از اشیای تیپ~$\sigma$ به اشیای تیپ~$\tau$ هستند.
برای نمونه، شاید بخواهیم تیپی~$\Omega$ از ارزش‌های صدق، «صادق» و
«کاذب»، و تیپی~$\Nat$ از اعداد طبیعی داشته باشیم. در این صورت، می‌توان
اشیای تیپ $\Nat \to \Omega$ را رابطه‌های یک‌موضعی یا زیرمجموعه‌های
$\Nat$ دانست؛ اشیای تیپ $\Nat \to \Nat$ تابع‌هایی از اعداد طبیعی به
اعداد طبیعی‌اند؛ و اشیای تیپ $(\Nat \to \Nat) \to \Nat$ «تابع‌تابع»ها،
یعنی تابع‌های تیپ بالاتری هستند که تابع‌ها را می‌گیرند و عدد پس می‌دهند.

همانند منطق مرتبهٔ دوم، می‌توان منطق مرتبهٔ بالاتر را نوعی منطق چندنوعی
دانست که برای هر تیپ از اشیای مورد نظرمان نوعی دارد. اما معمولاً روشن‌تر
است که نحو منطق تیپ‌های بالاتر را از پایه تعریف کنیم. برای نمونه،
می‌توانیم مجموعه‌ای از تیپ‌های متناهی را به صورت استقرایی چنین تعریف
کنیم:
\begin{enumerate}
\item $\Nat$ تیپی متناهی است.
\item اگر $\sigma$ و~$\tau$ تیپ‌های متناهی باشند، آنگاه $\sigma
  \to \tau$ نیز تیپی متناهی است.
\item اگر $\sigma$ و~$\tau$ تیپ‌های متناهی باشند، $\sigma \times
  \tau$ نیز تیپی متناهی است.
\end{enumerate}
به‌طور شهودی، $\Nat$ تیپ اعداد طبیعی، $\sigma \to \tau$ تیپ تابع‌ها
از~$\sigma$ به~$\tau$، و $\sigma \times \tau$ تیپ زوج‌هایی از اشیا را
نشان می‌دهد که یکی از~$\sigma$ و دیگری از~$\tau$ است. سپس می‌توانیم
مجموعه‌ای از ترم‌ها را به صورت استقرایی چنین تعریف کنیم:
\begin{enumerate}
\item برای هر تیپ~$\sigma$، ذخیره‌ای از متغیرهای $x$، $y$، $z$، \dots
  از تیپ~$\sigma$ وجود دارد.
\item $\Obj 0$ ترمی از تیپ~$\Nat$ است.
\item $\Obj S$ (جانشین) ترمی از تیپ~$\Nat \to \Nat$ است.
\item اگر $s$ ترمی از تیپ~$\sigma$ و $t$ ترمی از تیپ
  $\Nat \to (\sigma \to \sigma)$ باشد، آنگاه $\Obj{R}_{st}$ ترمی از
  تیپ $\Nat \to \sigma$ است.
\item اگر $s$ ترمی از تیپ~$\tau \to \sigma$ و $t$ ترمی از تیپ~$\tau$
  باشد، آنگاه $s(t)$ ترمی از تیپ~$\sigma$ است.
\item اگر $s$ ترمی از تیپ~$\sigma$ و $x$ متغیری از تیپ~$\tau$ باشد،
  آنگاه $\lambd[x][s]$ ترمی از تیپ $\tau \to \sigma$ است.
\item اگر $s$ ترمی از تیپ~$\sigma$ و $t$ ترمی از تیپ~$\tau$ باشد،
  آنگاه $\tuple{s, t}$ ترمی از تیپ $\sigma \times \tau$ است.
\item اگر $s$ ترمی از تیپ~$\sigma \times \tau$ باشد، آنگاه $p_1(s)$
  ترمی از تیپ~$\sigma$ و $p_2(s)$ ترمی از تیپ~$\tau$ است.
\end{enumerate}
به‌طور شهودی، $\Obj{R}_{st}$ تابعی را نشان می‌دهد که به صورت بازگشتی
زیر تعریف می‌شود:
\begin{align*}
\Obj{R}_{st}(0) & = s \\
\Obj{R}_{st}(x+1) & = t(x, R_{st}(x)),
\end{align*}
$\tuple{s, t}$ زوجی را نشان می‌دهد که مؤلفهٔ نخست آن~$s$ و مؤلفهٔ دوم
آن~$t$ است، و $p_1(s)$ و~$p_2(s)$ عنصر نخست و دوم («تصویرهای»)~$s$ را
نشان می‌دهند. سرانجام، $\lambd[x][s]$ تابع~$f$ را نشان می‌دهد که با
\[
f(x) = s
\]
برای هر~$x$ از تیپ~$\tau$ تعریف می‌شود؛ پس بند (6) صورتی از فراگیری به
ما می‌دهد و امکان می‌دهد تابع‌ها را با استفاده از ترم‌ها تعریف کنیم.
!!^{formula}s از عبارت‌های !!{identity}~$\eq[s][t]$ میان ترم‌های هم‌تیپ،
رابط‌های معمول گزاره‌ای، و کمیت‌گذاری تیپ بالاتر ساخته می‌شوند. سپس
می‌توان اصول موضوع دستگاه را معادلات پایهٔ حاکم بر ترم‌های تعریف‌شده در
بالا، همراه با قواعد معمول منطق دارای کمیت‌گذارها و !!{identity}، گرفت.

اگر دستگاه تیپ متناهی را با تیپی~$\Omega$ از ارزش‌های صدق گسترش دهیم،
باید اصول موضوعی را نیز بگنجانیم که بر کاربرد آن حاکم‌اند. در واقع، اگر
هوشمندانه عمل کنیم، می‌توانیم !!{formula}s مرکب را یکسره حذف کنیم و
به‌جای آن‌ها ترم‌هایی از تیپ~$\Omega$ بگذاریم!{} دستگاه اثبات نیز می‌تواند
متناسب با آن تغییر یابد. حاصل اساساً همان \emph{نظریهٔ سادهٔ تیپ‌ها}ست
که آلونزو چرچ در دههٔ ۱۹۳۰ عرضه کرد.

همانند منطق مرتبهٔ دوم، نسخه‌های متفاوتی از معناشناسی تیپ‌های بالاتر
وجود دارد که ممکن است بخواهیم به کار ببریم. در نسخهٔ تام، متغیرهای تیپ
$\sigma \to \tau$ روی مجموعهٔ \emph{همهٔ} تابع‌ها از اشیای تیپ~$\sigma$
به اشیای تیپ~$\tau$ دامنه می‌گیرند. چنان‌که می‌توان انتظار داشت، این
معناشناسی آن‌قدر قوی است که هیچ دستگاه !!{derivation} کامل و مؤثری
نمی‌پذیرد. اما می‌توان معناشناسی ضعیف‌تری را در نظر گرفت که در آن
!!a{structure} از مجموعه‌های عناصر~$T_\tau$ برای هر تیپ~$\tau$، همراه
با عمل‌های مناسب برای اعمال، تصویر و مانند آن، تشکیل می‌شود. اگر
جزئیات درست پرورده شوند، می‌توان قضایای تمامیت را برای انواع دستگاه‌های
!!{derivation} توصیف‌شده در بالا به دست آورد.

منطق تیپ‌های بالاتر جذاب است، زیرا چارچوبی فراهم می‌کند که در آن
می‌توان بخش بزرگی از ریاضیات را به‌شیوه‌ای طبیعی جای داد: با آغاز از
$\Nat$ می‌توان اعداد حقیقی، تابع‌های پیوسته، و مانند آن‌ها را تعریف
کرد. این منطق در زمینهٔ منطق شهودگرا نیز به‌ویژه جذاب است، زیرا تیپ‌ها
تعبیرهای «ساخت‌گرایانه» روشنی دارند. در واقع، می‌توان نسخه‌های
ساخت‌گرایانه‌ای از معناشناسی تیپ‌های بالاتر (مبتنی بر منطق شهودگرا، نه
منطق کلاسیک) پروراند که این تعبیرهای ساخت‌گرایانه را به‌خوبی روشن
می‌کنند و از بسیاری جهات از همتاهای کلاسیک خود جالب‌ترند.

\end{document}
